% Figure from MATH 493 notes, section 4. Compile with the notes preamble.
\begin{tikzpicture}[
    pt/.style={circle, draw, thick, minimum size=16pt, inner sep=0pt, fill=white, font=\small},
    arr/.style={-{Stealth[length=5pt]}, thick},
    ex/.style={font=\scriptsize}]
  % powers of 3 in U_7: 1, 3, 2, 6, 4, 5
  \begin{scope}[figB]
    \foreach \v [count=\k from 0] in {1,3,2,6,4,5} {
      \node[pt] (a\k) at ({90-60*\k}:1.35) {$\v$};
      \node[ex] at ({90-60*\k}:1.95) {$3^{\k}$};
    }
    \foreach \k [evaluate=\k as \j using {int(mod(\k+1,6))}] in {0,...,5}
      \draw[arr] (a\k) to[bend left=18] (a\j);
    \node[font=\small] at (0,-2.55) {$\langle 3 \rangle = U_7$, \ $\operatorname{ord}(3) = 6$};
  \end{scope}
  % powers of 2 in U_7: 1, 2, 4
  \begin{scope}[xshift=5.2cm]
    \begin{scope}[figR]
      \foreach \v [count=\k from 0] in {1,2,4} {
        \node[pt] (b\k) at ({90-120*\k}:1.0) {$\v$};
        \node[ex] at ({90-120*\k}:1.6) {$2^{\k}$};
      }
      \foreach \k [evaluate=\k as \j using {int(mod(\k+1,3))}] in {0,1,2}
        \draw[arr] (b\k) to[bend left=22] (b\j);
      \node[font=\small] at (0,-2.55) {$\langle 2 \rangle = \{1,2,4\}$, \ $\operatorname{ord}(2) = 3$};
    \end{scope}
    \foreach \v/\x in {3/-0.75, 5/0, 6/0.75}
      \node[pt, draw=black!40, text=black!55, thin] at (\x,-1.85+0.0) {$\v$};
  \end{scope}
\end{tikzpicture}
